\begin{tikzpicture}[font=\small,line cap=round,line join=round]\draw[otMuted](-1.4,0)--(2.6,0)node[right,font=\scriptsize]{$(d)$};\draw[geoTrait](0.750,1.500)circle(1.677);\draw[geoDonnee](1.500,0.000)--(0.000,3.000);\fill[otNavy](0.750,1.500)circle(1.3pt)node[right,geoLabel,font=\small]{$ O $};\fill[otNavy](0.000,0.000)circle(1.3pt)node[below left,geoLabel,font=\small]{$ A $};\fill[otNavy](1.500,0.000)circle(1.3pt)node[below right,geoLabel,font=\small]{$ B $};\fill[otNavy](0.000,3.000)circle(1.3pt)node[above left,geoLabel,font=\small]{$ C $};\geoSolution{\draw[otGreenDark,line width=.9pt](0.000,0.000)--(0.000,3.000);\draw[otGreenDark,line width=.6pt](0.170,0.000)--(0.170,0.170)--(0.000,0.170);}\end{tikzpicture}\par\vspace{1mm}{\fontsize{8}{10}\selectfont\color{otMuted}$[BC]$ est un diamètre du cercle de centre O passant par A.}